% Options for packages loaded elsewhere
\PassOptionsToPackage{unicode}{hyperref}
\PassOptionsToPackage{hyphens}{url}
%
\documentclass[
]{book}
\usepackage{amsmath,amssymb}
\usepackage{iftex}
\ifPDFTeX
  \usepackage[T1]{fontenc}
  \usepackage[utf8]{inputenc}
  \usepackage{textcomp} % provide euro and other symbols
\else % if luatex or xetex
  \usepackage{unicode-math} % this also loads fontspec
  \defaultfontfeatures{Scale=MatchLowercase}
  \defaultfontfeatures[\rmfamily]{Ligatures=TeX,Scale=1}
\fi
\usepackage{lmodern}
\ifPDFTeX\else
  % xetex/luatex font selection
\fi
% Use upquote if available, for straight quotes in verbatim environments
\IfFileExists{upquote.sty}{\usepackage{upquote}}{}
\IfFileExists{microtype.sty}{% use microtype if available
  \usepackage[]{microtype}
  \UseMicrotypeSet[protrusion]{basicmath} % disable protrusion for tt fonts
}{}
\makeatletter
\@ifundefined{KOMAClassName}{% if non-KOMA class
  \IfFileExists{parskip.sty}{%
    \usepackage{parskip}
  }{% else
    \setlength{\parindent}{0pt}
    \setlength{\parskip}{6pt plus 2pt minus 1pt}}
}{% if KOMA class
  \KOMAoptions{parskip=half}}
\makeatother
\usepackage{xcolor}
\usepackage{graphicx}
\makeatletter
\def\maxwidth{\ifdim\Gin@nat@width>\linewidth\linewidth\else\Gin@nat@width\fi}
\def\maxheight{\ifdim\Gin@nat@height>\textheight\textheight\else\Gin@nat@height\fi}
\makeatother
% Scale images if necessary, so that they will not overflow the page
% margins by default, and it is still possible to overwrite the defaults
% using explicit options in \includegraphics[width, height, ...]{}
\setkeys{Gin}{width=\maxwidth,height=\maxheight,keepaspectratio}
% Set default figure placement to htbp
\makeatletter
\def\fps@figure{htbp}
\makeatother
\setlength{\emergencystretch}{3em} % prevent overfull lines
\providecommand{\tightlist}{%
  \setlength{\itemsep}{0pt}\setlength{\parskip}{0pt}}
\setcounter{secnumdepth}{-\maxdimen} % remove section numbering
% definitions for citeproc citations
\NewDocumentCommand\citeproctext{}{}
\NewDocumentCommand\citeproc{mm}{%
  \begingroup\def\citeproctext{#2}\cite{#1}\endgroup}
\makeatletter
 % allow citations to break across lines
 \let\@cite@ofmt\@firstofone
 % avoid brackets around text for \cite:
 \def\@biblabel#1{}
 \def\@cite#1#2{{#1\if@tempswa , #2\fi}}
\makeatother
\newlength{\cslhangindent}
\setlength{\cslhangindent}{1.5em}
\newlength{\csllabelwidth}
\setlength{\csllabelwidth}{3em}
\newenvironment{CSLReferences}[2] % #1 hanging-indent, #2 entry-spacing
 {\begin{list}{}{%
  \setlength{\itemindent}{0pt}
  \setlength{\leftmargin}{0pt}
  \setlength{\parsep}{0pt}
  % turn on hanging indent if param 1 is 1
  \ifodd #1
   \setlength{\leftmargin}{\cslhangindent}
   \setlength{\itemindent}{-1\cslhangindent}
  \fi
  % set entry spacing
  \setlength{\itemsep}{#2\baselineskip}}}
 {\end{list}}
\usepackage{calc}
\newcommand{\CSLBlock}[1]{\hfill\break\parbox[t]{\linewidth}{\strut\ignorespaces#1\strut}}
\newcommand{\CSLLeftMargin}[1]{\parbox[t]{\csllabelwidth}{\strut#1\strut}}
\newcommand{\CSLRightInline}[1]{\parbox[t]{\linewidth - \csllabelwidth}{\strut#1\strut}}
\newcommand{\CSLIndent}[1]{\hspace{\cslhangindent}#1}
\ifLuaTeX
  \usepackage{selnolig}  % disable illegal ligatures
\fi
\usepackage{bookmark}
\IfFileExists{xurl.sty}{\usepackage{xurl}}{} % add URL line breaks if available
\urlstyle{same}
\hypersetup{
  pdftitle={Manual de Ecologia Evolutiva, Biometria e Biogeografia},
  pdfauthor={Renan Maestri},
  hidelinks,
  pdfcreator={LaTeX via pandoc}}

\title{Manual de Ecologia Evolutiva, Biometria e Biogeografia}
\author{Renan Maestri}
\date{2024-07-11}

\begin{document}
\frontmatter
\maketitle

\mainmatter
\chapter*{Prefácio}\label{prefuxe1cio}
\addcontentsline{toc}{chapter}{Prefácio}

\begin{center}\includegraphics[width=1\linewidth]{images/Capa do Livro} \end{center}

Esse manual tem o objetivo de servir de apoio às disciplinas que leciono
na UFRGS. O manual contém material de apoio das disciplinas de
Introdução à Ecologia, Ecologia para Biotecnologia, Ecologia I
(evolutiva) e Ecologia II (macroecologia e biogeografia), Ecologia e
Evolução Morfológica, Fundamentos de Morfometria Geométrica,
Macroecologia, Macroevolução e Métodos Filogenéticos Comparativos e
Oficina de Ecologia Evolutiva. Ministro essas disciplinas na
Universidade Federal do Rio Grande do Sul (UFRGS) desde 2018, algumas
para a graduação e algumas para a pós-graduação, para diferentes cursos
de graduação e Programas de Pós-Graduação do Instituto de Biociências da
UFRGS. Algumas dessas disciplinas são ou foram desenvolvidas em
colaboração com os professores Gonçalo Ferraz e Heinrich Hasenack
(Ecologia I), Leandro Duarte e Luciane Crossetti (Ecologia II), Rodrigo
Fornel (Fundamentos de Morfometria Geométrica) e Leandro Duarte (Oficina
de Ecologia Evolutiva), a quem agradeço pela parceria. O livro contém
uma parte do material visto nessas disciplinas, com foco nas atividades
práticas.

O intuito de disponibilizar esse manual de \textbf{Ecologia Evolutiva,
Biometria e Biogeografia} é que ele sirva de material de apoio às
disciplinas que ministro. Tentei organizá-lo de modo unificado, para que
ele seja útil simultaneamente às diferentes disciplinas e níveis de
ensino.

Todos os exercícios práticos usam o programa R (R Core Team 2023). Há um
capítulo que introduz o R.

Agradeço aos alunos com quem tive o prazer de interagir e que me
incentivam a ser um professor um pouco melhor a cada dia, espero que
esse manual seja útil a vocês durante e após as disciplinas. Tenho uma
dívida impagável com as pessoas que fazem parte e que passaram pelo
\href{https://www.ufrgs.br/lema/}{Laboratório de Ecomorfologia e
Macroevolução} pela troca de experiências e ensinamentos que culminaram
na elaboração deste manual.

Se você tiver qualquer comentário ou contribuição ao livro, por favor me
escreva
(\href{mailto:renanmaestri@gmail.com}{\nolinkurl{renanmaestri@gmail.com}}).

\phantomsection\label{refs}
\begin{CSLReferences}{1}{0}
\bibitem[\citeproctext]{ref-R-base}
R Core Team. 2023. \emph{R: A Language and Environment for Statistical
Computing}. Vienna, Austria: R Foundation for Statistical Computing.
\url{https://www.R-project.org/}.

\end{CSLReferences}

\backmatter
\end{document}
